% *======================================================================*
%  Cactus Thorn template for ThornGuide documentation
%  Author: Ian Kelley
%  Date: Sun Jun 02, 2002
%  $Header$
%
%  Thorn documentation in the latex file doc/documentation.tex
%  will be included in ThornGuides built with the Cactus make system.
%  The scripts employed by the make system automatically include
%  pages about variables, parameters and scheduling parsed from the
%  relevant thorn CCL files.
%
%  This template contains guidelines which help to assure that your
%  documentation will be correctly added to ThornGuides. More
%  information is available in the Cactus UsersGuide.
%
%  Guidelines:
%   - Do not change anything before the line
%       % START CACTUS THORNGUIDE",
%     except for filling in the title, author, date, etc. fields.
%        - Each of these fields should only be on ONE line.
%        - Author names should be separated with a \\ or a comma.
%   - You can define your own macros, but they must appear after
%     the START CACTUS THORNGUIDE line, and must not redefine standard
%     latex commands.
%   - To avoid name clashes with other thorns, 'labels', 'citations',
%     'references', and 'image' names should conform to the following
%     convention:
%       ARRANGEMENT_THORN_LABEL
%     For example, an image wave.eps in the arrangement CactusWave and
%     thorn WaveToyC should be renamed to CactusWave_WaveToyC_wave.eps
%   - Graphics should only be included using the graphicx package.
%     More specifically, with the "\includegraphics" command.  Do
%     not specify any graphic file extensions in your .tex file. This
%     will allow us to create a PDF version of the ThornGuide
%     via pdflatex.
%   - References should be included with the latex "\bibitem" command.
%   - Use \begin{abstract}...\end{abstract} instead of \abstract{...}
%   - Do not use \appendix, instead include any appendices you need as
%     standard sections.
%   - For the benefit of our Perl scripts, and for future extensions,
%     please use simple latex.
%
% *======================================================================*
%
% Example of including a graphic image:
%    \begin{figure}[ht]
% 	\begin{center}
%    	   \includegraphics[width=6cm]{MyArrangement_MyThorn_MyFigure}
% 	\end{center}
% 	\caption{Illustration of this and that}
% 	\label{MyArrangement_MyThorn_MyLabel}
%    \end{figure}
%
% Example of using a label:
%   \label{MyArrangement_MyThorn_MyLabel}
%
% Example of a citation:
%    \cite{MyArrangement_MyThorn_Author99}
%
% Example of including a reference
%   \bibitem{MyArrangement_MyThorn_Author99}
%   {J. Author, {\em The Title of the Book, Journal, or periodical}, 1 (1999),
%   1--16. {\tt http://www.nowhere.com/}}
%
% *======================================================================*

% If you are using CVS use this line to give version information
% $Header$

\documentclass{article}

% Use the Cactus ThornGuide style file
% (Automatically used from Cactus distribution, if you have a
%  thorn without the Cactus Flesh download this from the Cactus
%  homepage at www.cactuscode.org)
\usepackage{../../../../doc/latex/cactus}

\begin{document}

% The author of the documentation
\author{Roland Haas \textless rhaas@illinois.edu\textgreater \\ Erik Schnetter \textless eschnetter@perimeterinstitute.ca\textgreater}

% The title of the document (not necessarily the name of the Thorn)
\title{Backtrace}

% the date your document was last changed, if your document is in CVS,
% please use:
%    \date{$ $Date$ $}
% when using git instead record the commit ID:
%    \date{\gitrevision{<path-to-your-.git-directory>}}
\date{December 18 2024}

\maketitle

% Do not delete next line
% START CACTUS THORNGUIDE

% Add all definitions used in this documentation here
%   \def\mydef etc

% Add an abstract for this thorn's documentation
\begin{abstract}
This thorn produces stack traces in files \texttt{backtrace.<rank>.txt} when
the code receives a deadly signal.
It is agnostic of the driver being used and can be used with both PUGH and
CarpetX. Carpet(Lib) already contains its own (identical) code.
\end{abstract}

% The following sections are suggestive only.
% Remove them or add your own.

\section{Introduction}

This thorn registers signal handlers with the operating system to catch errors
caused by illegal memory access, calls to \code{abort()}, and similar
issues. When a signal is received it tries to output the call stack of
functions that lead to the failure. If possible it will resolve addresses to
symbol names and line numbers.

On systems where address space randomization is used, the addresses may only
be meaningful for the current process, though offsets to symbols will (should)
remain meaningful and can be used with \texttt{gdb} to resolve to a line
number.

To resolve symbols the code should be compiled with \texttt{-g} and
\texttt{-rdynamic} options (or equivalent).

\section{Numerical Implementation}

As of this writing the thorn captures the following signals:

\begin{lstlisting}[language=C]
  signal(SIGQUIT, signal_handler);
  signal(SIGILL, signal_handler);
  signal(SIGABRT, signal_handler);
  signal(SIGFPE, signal_handler);
  signal(SIGBUS, signal_handler);
  signal(SIGSEGV, signal_handler);
\end{lstlisting}

\section{Using This Thorn}

To use this thorn, it must be added to the \code{ActiveThorns} list in a
parameter file. No parameters need (or can) be set.

\subsection{Basic Usage}

\begin{lstlisting}
ActiveThorns = "Backtrace"
\end{lstlisting}

Backtrace files are named \texttt{backtrace.<rank>.txt} and are created in the
directory pointed to by \code{IO::out\_dir}.

\subsection{Special Behaviour}

The thorn will capture the current MPI rank very early during the run (during
\code{CCTK\_STARTUP}) and again once the driver has become active (during
\code{CCTK\_BASEGRID}) to report rank numbers and name files.

It will create the directory pointed to by \code{IO::out\_dir} if a backtrace
is requested and the directory does not yet exist.

\subsection{Interaction With Other Thorns}

This thorn uses the \code{out\_dir} parameter from thorn \texttt{IOUtil} to
decide which directory to put backtrace file in.

\subsection{Example output}
\begin{verbatim}
Backtrace from rank 0 pid 185787:
1. Backtrace::signal_handler(int)   [../../../../exe/cactus_backtrace(_ZN9Backtrace14signal_handlerEi+0xc4) [0x555e1fb92254]]
2. /lib/x86_64-linux-gnu/libc.so.6(+0x3fce0) [0x7f52b1049ce0]
3. /lib/x86_64-linux-gnu/libc.so.6(kill+0x7) [0x7f52b1049ef7]
4. ../../../../exe/cactus_backtrace(Backtrace_BacktraceTest+0x34) [0x555e1fb91f24]
5. ../../../../exe/cactus_backtrace(CCTK_CallFunction+0xd5) [0x555e2067c4b5]
6. ../../../../exe/cactus_backtrace(+0x103c992) [0x555e2067c992]
7. ../../../../exe/cactus_backtrace(+0x10433ad) [0x555e206833ad]
8. ../../../../exe/cactus_backtrace(CCTKi_DoScheduleTraverse+0x6b) [0x555e2068356b]
9. ../../../../exe/cactus_backtrace(+0x103c711) [0x555e2067c711]
10. ../../../../exe/cactus_backtrace(CCTK_ScheduleTraverse+0xd2) [0x555e2067e5a2]
11. ../../../../exe/cactus_backtrace(CCTKi_ScheduleTraverseGHExtensions+0x3c) [0x555e20671d3c]
12. ../../../../exe/cactus_backtrace(CactusDefaultInitialise+0xb9) [0x555e2066dc59]
13. ../../../../exe/cactus_backtrace(main+0x35) [0x555e1fa0dea5]
14. /lib/x86_64-linux-gnu/libc.so.6(+0x29d68) [0x7f52b1033d68]
15. /lib/x86_64-linux-gnu/libc.so.6(__libc_start_main+0x85) [0x7f52b1033e25]
16. ../../../../exe/cactus_backtrace(_start+0x21) [0x555e1fb871d1]

The hexadecimal addresses in this backtrace can also be interpreted
with a debugger (e.g. gdb), or with the 'addr2line' (or 'gaddr2line')
command line tool: 'addr2line -e cactus_sim <address>'.
\end{verbatim}

\subsection{Acknowledgements}

This thorn is almost entirely identical to \texttt{CarpetLib}'s implementation
found in git hash 99c01a802.

% Do not delete next line
% END CACTUS THORNGUIDE

\end{document}
